% =========================================================
% Well-Ordering and Equivalence
% =========================================================
\noindent\begin{tcolorbox}[colback=propbox, colframe=propborder, arc=2pt,
  left=6pt, right=6pt, top=4pt, bottom=4pt,
  title={\small\textbf{Well-Ordering Principle}},
  fonttitle=\small\bfseries]
Every nonempty subset of $\mathbb{N}$ has a least element.
\end{tcolorbox}
\label{sec:wellordering}

\begin{remark}[The minimal counterexample argument]
Well-ordering drives a proof strategy that is the structural dual of
induction: the \textbf{minimal counterexample argument}.

Instead of climbing up from a base case, you assume the claim fails
somewhere, take the smallest failure, and derive a contradiction from
its minimality. The logic:

\begin{enumerate}[itemsep=2pt]
  \item Suppose $P(n)$ fails for some $n$. Let
    $S = \{n \in \mathbb{N} : \neg P(n)\}$. By assumption $S \neq \emptyset$.
  \item By well-ordering, $S$ has a least element $m$.
  \item Derive a contradiction: either show $P(m)$ holds (contradicting
    $m \in S$), or find some $k < m$ with $\neg P(k)$ (contradicting
    the minimality of $m$).
\end{enumerate}

The key tool is minimality: the fact that $m$ is the \emph{smallest}
failure means all smaller values are successes, which is a powerful
assumption. It is equivalent to the strong inductive hypothesis.
\end{remark}

\begin{proposition}[Induction and well-ordering are equivalent]
Over the axioms of $\mathbb{N}$ without P5, the following three
principles are logically equivalent:
\begin{enumerate}[itemsep=2pt]
  \item The (weak) induction principle.
  \item The well-ordering principle.
  \item The strong induction principle.
\end{enumerate}
\end{proposition}

\begin{remark}[Proof sketch of the equivalence]
\textbf{Induction $\Rightarrow$ Well-ordering.}
Suppose $S \subseteq \mathbb{N}$ has no least element.
Let $P(n)$: ``$n \notin S$.''
Then $P(0)$ holds (otherwise $0 \in S$ would be the least element).
If $P(k)$ holds for all $k \leq n$, then $n+1 \notin S$ (otherwise
$n+1$ would be the least element of $S$). By induction, $P(n)$ holds
for all $n$, so $S = \emptyset$.

\medskip
\noindent\textbf{Well-ordering $\Rightarrow$ Induction.}
Suppose the base case $P(n_0)$ holds and the inductive step holds.
Let $S = \{n \geq n_0 : \neg P(n)\}$. If $S \neq \emptyset$, let $m$ be
its least element. Since $P(n_0)$ holds, $m \neq n_0$, so $m > n_0$
and $m - 1 \geq n_0$. By minimality of $m$, $P(m-1)$ holds.
But by the inductive step, $P(m-1) \Rightarrow P(m)$, contradicting
$m \in S$. Hence $S = \emptyset$.
\end{remark}

\begin{tcolorbox}[colback=gray!4, colframe=gray!35, arc=2pt,
  left=6pt, right=6pt, top=4pt, bottom=4pt,
  title={\small\textbf{Template: Minimal Counterexample}},
  fonttitle=\small\bfseries]
Suppose for contradiction that $P(n)$ fails for some $n \in \mathbb{N}$.

Let $S = \{n \in \mathbb{N} : \neg P(n)\}$.
By assumption $S \neq \emptyset$.
By well-ordering, $S$ has a least element $m$.

[\emph{Derive a contradiction using minimality of $m$: either show
$P(m)$ holds, or construct a smaller counterexample.}]

This contradicts [minimality of $m$ / the base case].
Therefore $P(n)$ holds for all $n \in \mathbb{N}$. \qed
\end{tcolorbox}

\begin{example}[Minimal counterexample: exponential bound]
\textbf{Claim.} For all $n \geq 1$, $2^n > n$.

\medskip
\noindent\textit{Proof (minimal counterexample).}
Suppose the claim fails. Let $S = \{n \geq 1 : 2^n \leq n\}$.
By assumption $S \neq \emptyset$; let $m$ be its least element.

Since $2^1 = 2 > 1$, we have $m \neq 1$, so $m \geq 2$.
Then $m - 1 \geq 1$ and $m - 1 < m$, so $m - 1 \notin S$
(by minimality of $m$). Thus $2^{m-1} > m - 1$.

Then:
\[
2^m = 2 \cdot 2^{m-1} > 2(m-1) = 2m - 2 \geq m
\]
(using $m \geq 2$ for the last inequality).

This contradicts $m \in S$ (which requires $2^m \leq m$).
Therefore $S = \emptyset$ and the claim holds for all $n \geq 1$. \qed

\medskip
\noindent
The minimal counterexample argument proceeds by taking the smallest
failure and showing it cannot be the smallest failure. The contradiction
is always of the same form: either the ``failure'' was not a failure,
or there is a smaller one.
\end{example}
